% Fixture: a book with chapters, numbering within chapters, running heads, author-year citations,
% a bibliography and a back-of-book index, with a PDF outline. Compile with: tectonic book-chapters.tex
\documentclass[11pt]{book}
\usepackage{amsmath,amssymb,amsthm}
\usepackage[paperwidth=6in,paperheight=8.5in,margin=0.8in]{geometry}
\usepackage[authoryear,round]{natbib}
\usepackage[bookmarks=true,hidelinks]{hyperref}
\newtheorem{theorem}{Theorem}[chapter]
\newtheorem{lemma}[theorem]{Lemma}
\newtheorem{proposition}[theorem]{Proposition}
\theoremstyle{definition}
\newtheorem{definition}[theorem]{Definition}
\newtheorem{exercise}{Exercise}[chapter]
\numberwithin{equation}{chapter}
\begin{document}
\frontmatter
\title{Notes on Linear Smoothers}
\author{A. Fixture}
\maketitle
\tableofcontents
\mainmatter
\chapter{Linear Smoothers}
\section{Least squares}
A \emph{linear smoother} is a method whose fitted values are a linear function of the response.
Least squares is the oldest example \citep{gauss1809}; a modern account is given by \citet{hastie2009}.
Let $X$ be the $N\times p$ input matrix. The fitted values are
\begin{equation}\label{eq:hat}
\hat y = X (X^T X)^{-1} X^T y = H y ,
\end{equation}
where $H$ is called the \emph{hat matrix}. The matrix in \eqref{eq:hat} is a projection.
\begin{definition}\label{def:df}
The \emph{degrees of freedom} of a linear smoother $S$ is $\operatorname{trace}(S)$.
\end{definition}
\begin{lemma}\label{lem:idem}
The hat matrix $H$ in \eqref{eq:hat} is symmetric and idempotent.
\end{lemma}
\begin{proof}
Direct computation gives $H^T = H$ and $H^2 = X (X^TX)^{-1} X^T X (X^TX)^{-1} X^T = H$.
\end{proof}
\begin{theorem}\label{thm:df}
The degrees of freedom of least squares, in the sense of Definition~\ref{def:df}, equals $p$.
\end{theorem}
\begin{proof}
By Lemma~\ref{lem:idem} and the cyclic property of the trace,
\begin{equation}\label{eq:trace}
\operatorname{trace}(H) = \operatorname{trace}\big((X^TX)^{-1} X^T X\big) = \operatorname{trace}(I_p) = p .
\end{equation}
This completes the argument.
\end{proof}
\section{Smoothing splines}
Smoothing splines minimize a penalized criterion \citep{reinsch1967, wahba1990}. Their fit is again linear in $y$, so Theorem~\ref{thm:df} generalizes through Definition~\ref{def:df}.
\begin{equation}\label{eq:spline}
\sum_{i=1}^{N} \{y_i - f(x_i)\}^2 + \lambda \int f''(t)^2\, dt .
\end{equation}
\begin{proposition}\label{prop:spline}
The minimizer of \eqref{eq:spline} is a natural cubic spline with knots at the unique values of $x_i$.
\end{proposition}
\begin{proof}
See \citet{wahba1990} for the classical argument; it uses \eqref{eq:spline} and an integration by parts.
\end{proof}
Figure~\ref{fig:smoother} compares the two smoothers of Section~\ref{sec:exercises}.
\begin{figure}[h]
\centering
\fbox{\rule{0pt}{0.8in}\hspace{2.6in}}
\caption{A smoothing spline and a least squares line fitted to the same data.}
\label{fig:smoother}
\end{figure}
\section{Exercises}\label{sec:exercises}
\begin{exercise}\label{ex:trace}
Show that the trace in \eqref{eq:trace} equals the rank of $X$.
\end{exercise}
\begin{exercise}
Use Proposition~\ref{prop:spline} to compute the degrees of freedom of a smoothing spline.
\end{exercise}

\chapter{Kernel Methods}
\section{Local averaging}
Kernel smoothers weight nearby observations \citep{nadaraya1964}. The Nadaraya--Watson estimate is
\begin{equation}\label{eq:nw}
\hat f(x_0) = \frac{\sum_{i=1}^{N} K_\lambda(x_0, x_i)\, y_i}{\sum_{i=1}^{N} K_\lambda(x_0, x_i)} ,
\end{equation}
where $K_\lambda$ is a kernel with bandwidth $\lambda$. Like \eqref{eq:hat}, it is a linear smoother.
\begin{theorem}\label{thm:nw}
The Nadaraya--Watson estimator \eqref{eq:nw} is a linear smoother whose degrees of freedom, as in Theorem~\ref{thm:df}, depend on $\lambda$.
\end{theorem}
\begin{proof}
The weights in \eqref{eq:nw} do not depend on $y$, so by Definition~\ref{def:df} the claim follows from Lemma~\ref{lem:idem}.
\end{proof}
Exercise~\ref{ex:trace} in Chapter~1 is useful here. \citet{hastie2009} discuss bandwidth selection at length.

\backmatter
\bibliographystyle{plainnat}
\begin{thebibliography}{6}
\bibitem[Gauss(1809)]{gauss1809} Gauss, C. F. (1809). Theoria Motus Corporum Coelestium. Perthes et Besser, Hamburg.
\bibitem[Hastie et~al.(2009)Hastie, Tibshirani, and Friedman]{hastie2009} Hastie, T., Tibshirani, R. and Friedman, J. (2009). The Elements of Statistical Learning. Springer, New York, second edition.
\bibitem[Nadaraya(1964)]{nadaraya1964} Nadaraya, E. A. (1964). On estimating regression. Theory of Probability and its Applications, 9:141--142.
\bibitem[Reinsch(1967)]{reinsch1967} Reinsch, C. (1967). Smoothing by spline functions. Numerische Mathematik, 10:177--183.
\bibitem[Wahba(1990)]{wahba1990} Wahba, G. (1990). Spline Models for Observational Data. SIAM, Philadelphia.
\end{thebibliography}
\begin{theindex}
\item degrees of freedom, 2, 5
\item hat matrix, 2
\item kernel smoother, 5
\item linear smoother, 1, 5
\subitem degrees of freedom, 2
\item Nadaraya--Watson estimate, 5
\item smoothing spline, 3--4
\item trace, 2
\end{theindex}
\end{document}
